\documentclass[10pt,a4paper]{amsart}
\usepackage[margin=19mm]{geometry}
\usepackage{amsmath,amssymb,mathtools}
\usepackage[scr=rsfs,cal=euler]{mathalfa}
\usepackage{xfrac}
\usepackage[unicode,hidelinks,bookmarks=false]{hyperref}
\newcommand{\ie}{i.e.\ }
\newcommand{\cf}{cf.\ }
\newcommand{\resp}{respectively\ }
\newcommand{\Subsection}{\subsection}
\providecommand{\nobreakdash}{\nobreak\hbox{-}}
\setlength{\emergencystretch}{2em}
\multlinegap0pt
\usepackage{amssymb}
\newtheorem{lemma}{Lemma}
\title{Sharp large deviation estimates for Gaussian extrema}
\author{Jos\'e M. Zapata}
\date{Source: arXiv:2512.18297v1 (CC BY 4.0)}
\begin{document}
\maketitle

\noindent\emph{Source and licence.} Sharp large deviation estimates for Gaussian extrema. Version arXiv:2512.18297v1 (CC BY 4.0), distributed by arXiv under the Creative Commons Attribution 4.0 International licence, http://creativecommons.org/licenses/by/4.0/, as recorded in the arXiv licence metadata for this exact version and shown on the abstract page https://arxiv.org/abs/2512.18297v1. Full licence notice: \texttt{licenses/arxiv-2512.18297-CC-BY-4.0.txt}.\par\smallskip

\noindent\textit{Editorial scope.} Counted unit: the article's lemma lem:density (source0.tex lines 292--301). The article's complete original proof of it is delivered with it (source0.tex lines 303--375, one \texttt{proof} environment of 73 lines). No mathematics is written, repaired or re-derived: the statement and the proof are the article's own lines, in the article's own notation, and no retained line is substituted or re-flowed.  Retained uncounted as the source-local context the proof needs: lines 217--228; lines 219--222.  Nothing was omitted from any retained span, so the package carries no omission record.  No retained line contains a citation, so the wrapper carries no bibliography and no bibliography entry.  No retained line contains a figure environment, an image-inclusion command or drawing code, so no image is delivered and the package declares no asset dependency.  Editorial formatting only, in addition: an AMS \texttt{amsart} preamble that restates the article's own declarations and macros used by the retained lines, namely the environments \texttt{lemma} on separate counters, together with the standard packages those lines need.  The retained environments print in this document as Lemma 1 in order of appearance, whereas the article prints the same items with the numbering of its own full document.  The complete native source of the article as distributed in the arXiv e-print for this exact version is bundled unchanged as \texttt{sources/191374/source0.tex}.
\begin{lemma}\label{lem:ineqNormal}
There exists $t_0$ such that
\begin{equation}\label{eq:normalI}
0 \le -\log \Phi(t) \le \frac{2\,\phi(t)}{t},
\qquad \mbox{ for all }t>t_0.
\end{equation}
Moreover, 
\begin{equation}\label{eq:normalII}
1-\Phi(t)^n \le n\,\frac{\phi(t)}{t},
\qquad\mbox{ for all }t>0\mbox{ and }n\in\mathbb{N}.
\end{equation}
\end{lemma}

\begin{equation}
0 \le -\log \Phi(t) \le \frac{2\,\phi(t)}{t},
\qquad \mbox{ for all }t>t_0.
\end{equation}

\begin{lemma}\label{lem:density}
Let $Z_n$ be defined as above and let $f_n$ denote its density.
Then, for any fixed $M>0$,  
\[
\lim_{n\to\infty}\frac{1}{\log n}\log f_n(x)
=
-I(x),
\]
 uniformly on $x\in[0,M]$.
\end{lemma}

\begin{proof}
Since $a_n^2=2\log n$, we may write
\[
\frac{1}{\log n}\log f_n(x)
=
\frac{2}{a_n^2}\log\!\left(\frac{n a_n}{2}\right)
+
\frac{2(n-1)}{a_n^2}\log\Phi\!\bigl(t_n(x)\bigr)
+
\frac{2}{a_n^2}\log\phi\!\bigl(t_n(x)\bigr).
\]
Denote by $T_i(n,x)$, $i=1,2,3$, the i-th term in the sum above. We compute the limit of each of these terms and verify that the convergence is  uniform in  $[0,M]$.

First,
\[
\lim_{n\to\infty}T_1(n,x)=
1+\lim_{n\to\infty}\frac{\log a_n-\log 2}{\log n}=1.
\]
This convergence is obviously uniform since the
term does not depend on $x$.

Second, applying \eqref{eq:normalI} in Lemma \ref{lem:ineqNormal} taking into account that $t_n(x)\ge t_n(0)=b_n\to \infty$, we have
\begin{equation}\label{eq:T2}
|T_2(n,x)|\le-\frac{2(n-1)}{a_n^2}\log\Phi(b_n)\le 
\frac{4(n-1)}{a_n^2}\frac{\phi(b_n)}{b_n}
\end{equation}
for $n$ large enough and all $x\ge 0$. 
Inspection shows that  $n\phi(b_n)/b_n\to 1$ as $n\to \infty$. Then, the upper bound in \eqref{eq:T2} converges to $0$ as $n\to\infty$. Since it does not depend on $x$ the convergence is  uniform in  $x$. 


Third, using that
\[
t_n(x)=a_n\left(1+\frac{x}{2}\right)-\frac{\log\left(\sqrt{2\pi}a_n\right)}{a_n},
\]
one gets
\begin{align*}
T_3(n,x)
&=
-\left(1+\frac{x}{2}\right)^2
+
\frac{2\left(1+\frac{x}{2}\right)\log(\sqrt{2\pi}a_n)}{a_n^2}  %\\&\quad
-\frac{\log(2\pi)}{a_n^2}
-
\frac{\log^2(\sqrt{2\pi}a_n)}{a_n^4}.
\end{align*}
Then, for $x\in[0,M]$, we have
\[
\left|T_3(n,x)+\left(1+\frac{x}{2}\right)^2\right|
\le 
2\left(1+\frac{M}{2}\right)\frac{\log(\sqrt{2\pi}a_n)}{a_n^2}+\frac{\log(2\pi)}{a_n^2}
+
\frac{\log^2(\sqrt{2\pi}a_n)}{a_n^4}.
\]
The right hand side of the above inequality does not depend on $x$ and converges to $0$ as $n\to\infty$. Hence,
\[
T_3(n,x)
\longrightarrow
-\left(1+\frac{x}{2}\right)^2,\quad\mbox{ as }
n\to\infty
\]
uniformly for $x\in[0,M]$.


Combining the previous limits yields, for every $x\in[0,M]$,
\[
\lim_{n\to\infty}\frac{1}{\log n}\log f_n(x)
=
1-\left(1+\frac{x}{2}\right)^2
=
-x-\frac{x^2}{4}=-I(x).
\]
The convergence is  uniform in  the interval $[0,M]$. The proof is complete.
\end{proof}

\end{document}
